\documentclass [tikz] {standalone}
\input{header.htex}


\begin {document}
\begin{tikzpicture}
	
	\begin{axis}[
		width = 12cm,
		height = 7cm,
		xshift = 0cm,
		xmin = 0,
		xmax = 1.0,
		ymin = -0.3,
		ymax = 1.5,
		xlabel = {Нормированная частота, $f/f_д$},
		ylabel = {Амплитуда}, 
		xtick={0, 0.25 , 0.5, 0.75, 1.0},
		xticklabels={0, 0.125 , 0.25, 0.375, 0.5},
		grid = major,
		minor tick num=3]
		\node [fill = white, above, rotate = 0]  at (axis cs: 0.45, 0.5){$\alpha$};
		\node [fill = white, above, rotate = 0]  at (axis cs: 0.55, 0.5){$\alpha$};
		
		\addplot + [line width = 1.5 pt,  draw=black, mark = none,] table {../text data/filter_change.txt};
        \addlegendentry{$\text{АЧХ HBF-фильтра}$};
        
        \draw[-] (axis cs: 0, 0.95) -- (axis cs: 0.46, 0.95);
        \draw[-] (axis cs: 0, 1.05) -- (axis cs: 0.46, 1.05);
        \draw[->] (axis cs: 0.08, 1.05) -- (axis cs: 0.08, 1.25);
        \draw[->] (axis cs: 0.08, 0.95) -- (axis cs: 0.08, 0.75);
        \node [fill = white, above, rotate = 0]  at (axis cs: 0.36, 1.07){Пульсации в полосе пропускания};
        
        \draw[-] (axis cs: 0.55, 0) -- (axis cs: 1, 0);
        \draw[-] (axis cs: 0.55, 0.05) -- (axis cs: 1, 0.05);
        \draw[->] (axis cs: 0.625, 0.05) -- (axis cs: 0.625, 0.255);
        \draw[->] (axis cs: 0.625, 0) -- (axis cs: 0.625, -0.21);
%        \node [fill = white, above, rotate = 0]  at (axis cs: 0.7, 0.06){Пульсации в полосе \\ задерживания};
		\node[
		fill=white,
		above,
		rotate=0,
		align=center
		]
		at (axis cs: 0.8, 0.06)
		{Пульсации в полосе\\задерживания};
        
        \draw[-] (axis cs: 0.4, 1.1) -- (axis cs: 0.4, -0.1);
        \draw[-] (axis cs: 0.6, 1.1) -- (axis cs: 0.6, -0.1);
        
        \draw[<->] (axis cs: 0.4, 0.5) -- (axis cs: 0.5, 0.5);
        \draw[<->] (axis cs: 0.5, 0.5) -- (axis cs: 0.6, 0.5);\
        \draw[->] (axis cs: 0.4, 0.6) -- (axis cs: 0.30, 0.6);
        \draw[->] (axis cs: 0.6, 0.6) -- (axis cs: 0.70, 0.6);
        
    \end{axis}
			
\end{tikzpicture}
\end {document}